В ходе работы были рассмотрены два классических алгоритма сортировки: быстрая сортировка
и сортировка слиянием. Показано, что оба алгоритма обладают средней временной сложностью
$O(n \log n)$, однако различаются по поведению в наихудшем случае и объёму используемой
дополнительной памяти (см.~таблицу~\ref{tab:comparison}).

Анализ рекуррентных соотношений~\eqref{eq:quicksort-recurrence}
и~\eqref{eq:mergesort-recurrence} подтверждает теоретические оценки, приведённые
в~\cite{cormen2009algorithms,knuth1997art}. Практические реализации на C++ и Python
демонстрируют основные структурные различия двух подходов.

Выбор между алгоритмами определяется условиями задачи: быстрая сортировка предпочтительна
при ограниченной памяти и случайных данных, тогда как сортировка слиянием гарантирует
$O(n \log n)$ в наихудшем случае и является устойчивой.
